% =============================================================
%  ch11_installation.tex  --  ROS 2 installieren (Schritt für Schritt)
% =============================================================

\chapter{ROS\,2 installieren -- Schritt für Schritt}

\begin{lstlisting}[style=bash,caption={Locale + ROS\,2-Apt-Quelle (Ubuntu 24.04)}]
# 1) Locale auf UTF-8
sudo apt update && sudo apt install -y locales
sudo locale-gen en_US en_US.UTF-8
sudo update-locale LC_ALL=en_US.UTF-8 LANG=en_US.UTF-8

# 2) Universe-Repo + Schluessel + ROS2-Quelle
sudo apt install -y software-properties-common curl
sudo add-apt-repository universe
sudo curl -sSL https://raw.githubusercontent.com/ros/rosdistro/master/ros.key \
  -o /usr/share/keyrings/ros-archive-keyring.gpg
echo "deb [arch=$(dpkg --print-architecture) \
  signed-by=/usr/share/keyrings/ros-archive-keyring.gpg] \
  http://packages.ros.org/ros2/ubuntu \
  $(. /etc/os-release && echo $UBUNTU_CODENAME) main" \
  | sudo tee /etc/apt/sources.list.d/ros2.list > /dev/null
sudo apt update
\end{lstlisting}

\begin{lstlisting}[style=bash,caption={ROS\,2 Jazzy + Entwicklungswerkzeuge}]
# 3) ROS2 Desktop (inkl. RViz2, Demos)
sudo apt install -y ros-jazzy-desktop

# 4) Build-Werkzeuge
sudo apt install -y ros-dev-tools \
  python3-colcon-common-extensions python3-rosdep
sudo rosdep init && rosdep update

# 5) In jeder Shell verfuegbar machen
echo "source /opt/ros/jazzy/setup.bash" >> ~/.bashrc && source ~/.bashrc

# 6) Funktionstest (zwei Terminals)
ros2 run demo_nodes_cpp talker      # Terminal A
ros2 run demo_nodes_py listener     # Terminal B
\end{lstlisting}

\section{Pakete, die du für die Projekte brauchst}
\begin{longtable}{@{}p{6.2cm}p{8.6cm}@{}}
\toprule
\textbf{Paket (apt: \code{ros-jazzy-...})} & \textbf{Zweck} \\
\midrule
\code{ros-gz} (Gazebo Harmonic Bridge)      & Simulation $\leftrightarrow$ ROS\,2 \\
\code{ros2-control}, \code{ros2-controllers}& Hardware-Abstraktion, \code{joint\_trajectory\_controller} \\
\code{gz-ros2-control}                      & ros2\_control-Plugin für Gazebo \\
\code{moveit}                               & Bewegungsplanung, IK, Greifen (Projekt~4) \\
\code{navigation2}, \code{nav2-bringup}     & Autonome Navigation (Projekt~3) \\
\code{slam-toolbox}                         & SLAM/Kartierung (Projekt~3) \\
\code{robot-localization}                   & EKF/UKF-Sensorfusion (Odometrie) \\
\code{image-pipeline}                       & \code{camera\_calibration}, \code{stereo\_image\_proc} \\
\code{v4l2-camera} bzw.\ \code{usb-cam}     & Treiber für die USB-Stereokamera \\
\code{vision-opencv} (\code{cv\_bridge})    & OpenCV $\leftrightarrow$ ROS\,2-Bildnachrichten \\
\code{xacro}, \code{robot-state-publisher}  & URDF/Xacro, TF-Veröffentlichung \\
\code{rqt}, \code{rqt-graph}, \code{rqt-image-view} & Introspektion/Visualisierung \\
\bottomrule
\end{longtable}
\noindent Installation gebündelt, z.\,B.\
\code{sudo apt install ros-jazzy-ros-gz ros-jazzy-ros2-control ros-jazzy-moveit \ldots}\
Für den ESP32 zusätzlich der \textbf{micro-ROS Agent} (separat aus dem
micro-ROS-Workspace gebaut).
